\chapter{The space-form layer: definitions, flow, admissibility}
\label{chap:spaceform_core}

% Chapter for the K-generic space-form engine:
%   Gluck/SpaceForm/Defs.lean          — Realizes, SpaceFormFourVertex, spaceFormSpeed,
%                                        centeredRadius and its algebra
%   Gluck/SpaceForm/Flow.lean          — truncatedSpeed/Field, Picard–Lindelöf flow
%                                        spaceFormFlow, endpoint map
%   Gluck/SpaceForm/Admissible.lean    — Grönwall L¹ drive, invariant admissible domain
%   Gluck/SpaceForm/ArcAlgebra.lean    — constant-curvature arcs, arc-endpoint map,
%                                        four-arc step error map, single-arc transport
%   Gluck/SpaceForm/StepReparam.lean   — quarter-step and four-arc step-model transport
%   Gluck/SpaceForm/Margins.lean       — arcMargins package, stepModel_margins
%   Gluck/SpaceForm/FirstVariation.lean — stepError_expansion, stepError_coeff_ne_zero
%
% This is the spherical flow engine of chap:sphere transported over the
% ambient curvature K ∈ {+1, −1, 0}: K = +1 recovers S²
% verbatim, K = −1 is the Poincaré-disk model of H², and K = 0 is the flat
% Euclidean plane in the conformal gauge λ = 2. The model-agnostic step
% machinery (stepCurvature, exists_step_L1_reparam, the FirstVariation frame)
% is REUSED from Gluck/Sphere/* via imports; only the metric-carrying layers
% are re-derived with the K parameter.

\section{Overview}

This chapter generalizes the spherical engine of \cref{chap:sphere} over the
\emph{ambient curvature} $K \in \{+1,-1,0\}$ (the constant sectional
curvature of the space form). All three constant-curvature models live in the
open unit disk $\{|z|<1\}\subseteq\mathbb C$ in the tangent-angle gauge: the
sphere $S^2$ ($K=+1$, round metric $4/(1+|z|^2)^2$), the
hyperbolic plane $H^2$ ($K=-1$, Poincar\'e metric $4/(1-|z|^2)^2$),
and the Euclidean plane $E^2$ ($K=0$, flat conformal
metric $\lambda=2$, so that the disk-gauge curvature is half the Euclidean
one). The single unifying object is the \emph{space-form
gauge speed}
\[
  q_{K,\kappa}(\theta,z)
  \;=\;\frac{1+K\|z\|^2}
        {2\bigl(\kappa(\theta)-K\langle z, ie^{i\theta}\rangle\bigr)},
\]
the algebraic solution of the conformal geodesic-curvature relation for the
speed $\|\gamma'\|$ in the gauge $\varphi(\theta)=\theta$. At $K=+1$ this
is the landed \texttt{Gluck.sphericalSpeed} (\cref{def:spherical_speed}); at
$K=-1$ it is the hyperbolic speed. Both the metric factor
$1+K\|z\|^2$ and the inner-product term $-K\langle
\cdot\rangle$ flip with $K$; that the two sign flips are a
\emph{single} parameter is verified once, in the model-circle anchor
\cref{lem:spaceform_speed_circle}.

The chapter is a $K$-generic transport of the spherical files
\texttt{Gluck/Sphere/\{Defs, Flow, Admissible, ArcAlgebra, StepReparam,
Margins, FirstVariation\}}, and it is honest about what is transported versus
reused: the metric-free step machinery — the canonical four-arc
\texttt{stepCurvature} (\cref{def:four_arc_step_function}), the $L^1$
step reparametrization \cref{lem:step_L1_reparam}, and the FirstVariation
frame algebra — is pure step-function/measure theory with no $K$,
and is \emph{imported verbatim} from the \texttt{Gluck/Sphere} layer (per the
module imports of \texttt{StepReparam.lean} and \texttt{FirstVariation.lean}).
Only the metric-carrying layers are re-derived with the curvature parameter.

Three $K$-specific phenomena organize the transport:

\begin{itemize}
\item \emph{The position-dependent admissibility denominator.} Every estimate
  runs through the bracket
  $\kappa(\theta)-K\langle z, ie^{i\theta}\rangle$. For
  $|K|\le1$, Cauchy--Schwarz gives
  $|K\langle z, ie^{i\theta}\rangle|\le\|z\|$, so the bracket is
  positive on $\{\|z\|\le R<\kappa\}$ \emph{uniformly in} $K$
  (\cref{lem:spaceform_denom_pos}) — most of the flow layer needs only
  $|K|\le1$, not the sign. The truncated flow clamps this
  denominator at $\delta$ and the norm at $R$, exactly as on the sphere.
\item \emph{The confinement threshold.} Where the admissibility lower bound
  is \emph{established} per model, the sign enters, through the single
  uniform threshold $(1-K)/2$: for $K=+1$ it is $0$
  (positivity of $\kappa$ suffices), for $K=-1$ it is $1$ (the
  escape-velocity bound — geodesic circles in $H^2$ have curvature
  $\coth\rho>1$, and curves with $\kappa\le1$ never close), and for
  $K=0$ it is $1/2$ (the Euclidean circle of disk-gauge curvature
  $c$ has coordinate radius $1/(2c)$, which fits in the unit disk exactly
  when $c>1/2$). The four-vertex hypothesis
  (\cref{def:spaceform_four_vertex}) carries this threshold, and the
  centered radius clears the disk precisely above it
  (\cref{lem:centered_radius_spaceform_props}).
\item \emph{No absolute floor on the bracket scale.} The bracket value at the
  level-$c$ model circle is $B=c+K r^\ast=\sqrt{c^2+K}$
  (\cref{lem:centered_radius_spaceform_props}). For $S^2$ this is
  $\sqrt{1+c^2}\ge1$, an absolute lower bound; for $H^2$ it is
  $\sqrt{c^2-1}\to0$ as $c\to1^{+}$. The margin constants of
  \cref{lem:spaceform_step_model_margins} are therefore fractions of $B$
  ($c$-dependent), not of $1$ — the one place the sphere's
  absolute-constants shortcut breaks. Relatedly, the first-variation
  coefficient $\eta(K)=2K r^\ast/(c^2+K)$
  is positive for $K=+1$ but \emph{negative} for
  $K=-1$; the winding argument downstream consumes only
  $\eta\ne0$ (\cref{lem:spaceform_step_error_coeff_ne_zero}). At
  $K=0$ the coefficient \emph{vanishes} — the flat gauge field is
  position-independent, so the start-point degree of freedom is dead — and
  the flat member exits the endpoint-winding pipeline there, closing instead
  through the classical alignment winding
  (\cref{lem:reduction_justified_L1}; see \cref{chap:spaceform_converse}).
\end{itemize}

The declarations below live in \texttt{namespace Gluck.SpaceForm}. The
capstone consumers (endpoint winding, reconstruction, the positive and
mixed-sign space-form converses) are stated in \cref{chap:spaceform_converse}.

\section{Declarations}

\subsection{Definition layer}

\begin{definition}\leanok
[space-form realization]
  \label{def:spaceform_realizes}
  \lean{Gluck.SpaceForm.Realizes}
  \uses{def:realizes_spherical_curvature, def:realizes_curvature}
  A curve $\gamma:\mathbb R\to\mathbb C$ \emph{realizes the space-form curvature
  function} $\kappa$ at ambient curvature $K$ when it is $C^1$, regular
  ($\gamma'(t)\ne0$), confined to the open unit disk ($\|\gamma(t)\|<1$), and there is a
  differentiable tangent angle $\varphi$ with
  $\gamma'(t)=\|\gamma'(t)\|\,e^{i\varphi(t)}$ and
  \[
    \frac{1+K\|\gamma(t)\|^2}{2}\,\varphi'(t)
    \;=\;\bigl(\kappa(t)-K\langle \gamma(t),
      ie^{i\varphi(t)}\rangle\bigr)\,\|\gamma'(t)\|.
  \]
  At $K=+1$ this is \cref{def:realizes_spherical_curvature}; at
  $K=-1$ it is the Poincar\'e-disk statement that $\gamma$ has
  hyperbolic geodesic curvature $\kappa$ (up to the tangent-angle gauge).
\end{definition}

\begin{definition}\leanok
[space-form four-vertex hypothesis]
  \label{def:spaceform_four_vertex}
  \lean{Gluck.SpaceForm.SpaceFormFourVertex}
  \uses{def:is_curvature_function, def:four_vertex_condition,
        def:sphere_four_vertex}
  The $K$-generic admissibility hypothesis for the positive-stage
  converse: $\kappa$ is a curvature function
  (\cref{def:is_curvature_function}: continuous, $2\pi$-periodic, strictly
  positive), satisfies the value-separated four-vertex condition
  \cref{def:four_vertex_condition}, and — for the non-spherical members
  $K\le0$ — the uniform confinement floor
  $\kappa(\theta)>(1-K)/2$ for all $\theta$. At $K=+1$
  the last clause is vacuous and this is \cref{def:sphere_four_vertex} (with
  a trivially true extra conjunct); at $K=-1$ it is the
  escape-velocity bound $\kappa>1$, the threshold below which curves in
  $H^2$ cannot close; at $K=0$ it is the flat threshold
  $\kappa>1/2$, which makes the model circle of any level in the range of
  $\kappa$ fit inside the open unit disk
  (\cref{lem:centered_radius_spaceform_props}).
\end{definition}

\begin{definition}\leanok
[space-form gauge speed]
  \label{def:spaceform_speed}
  \lean{Gluck.SpaceForm.spaceFormSpeed}
  \uses{def:spherical_speed}
  The \emph{space-form gauge speed}
  \[
    q_{K,\kappa}(\theta,z)=
    \frac{1+K\|z\|^2}
         {2\,(\kappa(\theta)-K\langle z, ie^{i\theta}\rangle)},
  \]
  a junk-value total function; every lemma about it carries the admissibility
  hypotheses $\|z\|\le R<\kappa(\theta)$ (with $|K|\le1$) making the
  denominator positive. At $K=+1$ it is \cref{def:spherical_speed}.
\end{definition}

\begin{lemma}\leanok
[$K$-uniform denominator and speed positivity]
  \label{lem:spaceform_denom_pos}
  \lean{Gluck.SpaceForm.denom_pos}
  \uses{def:spaceform_speed, lem:sphere_denom_pos}
  For $|K|\le1$ and $\|z\|\le R<\kappa(\theta)$ the admissibility
  bracket $\kappa(\theta)-K\langle z, ie^{i\theta}\rangle$ is
  strictly positive; consequently (for $\|z\|<1$, where the numerator
  $1+K\|z\|^2>0$) the gauge speed is strictly positive. Recovers
  \cref{lem:sphere_denom_pos} at $K=1$ and the hyperbolic bound at
  $K=-1$ with one proof.
\end{lemma}

\begin{proof}
  \leanok
  Cauchy--Schwarz against the unit vector $ie^{i\theta}$:
  $|K\langle z, ie^{i\theta}\rangle|\le|K|\,\|z\|\le
  \|z\|\le R<\kappa(\theta)$, then \texttt{linarith}. Numerator positivity is
  $\|z\|^2<1$ and $K\ge-1$ (\texttt{nlinarith}).
\end{proof}

\begin{lemma}\leanok
[model geodesic-circle anchor (sign linchpin)]
  \label{lem:spaceform_speed_circle}
  \lean{Gluck.SpaceForm.spaceFormSpeed_circle}
  \uses{def:spaceform_speed, lem:spherical_speed_circle}
  The Euclidean circle of radius $r$ traversed in the tangent-angle gauge sits
  at $\gamma(\theta)=(-r)\,ie^{i\theta}$ and carries constant space-form curvature
  $\kappa=(1-K r^2)/(2r)$; its gauge speed comes out to exactly
  $r$. Recovers the sphere circle $\kappa_S=(1-r^2)/(2r)=\cot\rho$
  ($K=1$) and the hyperbolic circle
  $\kappa_H=(1+r^2)/(2r)=\coth\rho$ ($K=-1$). The opposite
  inner-product sign convention would give junk here — this lemma pins the
  two-signs-are-one-parameter convention for the whole layer.
\end{lemma}

\begin{proof}
  \leanok
  Direct computation: $\langle(-r)ie^{i\theta}, ie^{i\theta}\rangle=-r$ and
  $\|(-r)ie^{i\theta}\|=r$, so the denominator is
  $2(\kappa+K r)=(1+K r^2)/r$ and the quotient collapses
  to $r$ under $1+K r^2\ne0$ (\texttt{field\_simp}).
\end{proof}

\begin{definition}\leanok
[space-form centered radius]
  \label{def:centered_radius}
  \lean{Gluck.SpaceForm.centeredRadius}
  \uses{lem:centered_radius_pos}
  The \emph{space-form centered radius}, in the rationalized uniform form
  \[
    r^\ast(K,c)\;=\;
    \frac{1}{\sqrt{c^2+K}+c},
  \]
  the root in $(0,1)$ of $K r^2+2cr-1=0$, i.e.\ the Euclidean radius
  of the model geodesic circle of constant curvature $c$. The single closed
  form recovers the sphere radius $\sqrt{1+c^2}-c$ of
  \cref{lem:centered_radius_pos} ($K=1$), the hyperbolic radius
  $c-\sqrt{c^2-1}$ ($K=-1$, for $c>1$), and — unlike the legacy
  unrationalized form $K(\sqrt{c^2+K}-c)$, which
  vanishes identically at $K=0$ — the flat radius $1/(2c)$
  ($K=0$, for $c>1/2$). The two forms agree at
  $K=\pm1$ (\texttt{centeredRadius\_eq}).
\end{definition}

\begin{lemma}\leanok
[centered-radius algebra]
  \label{lem:centered_radius_spaceform_props}
  \lean{Gluck.SpaceForm.centeredRadius_solves,
        Gluck.SpaceForm.centeredRadius_mem_Ioo,
        Gluck.SpaceForm.centeredRadius_bracket,
        Gluck.SpaceForm.centeredRadius_facts,
        Gluck.SpaceForm.centeredRadius_lt_one_iff}
  \uses{def:centered_radius}
  For $K\in\{+1,-1,0\}$ and an admissible level $c$ (meaning
  $c>0$ when $K=1$, $c>1$ when $K=-1$, $c>1/2$ when
  $K=0$), the centered radius
  $r^\ast=r^\ast(K,c)$ satisfies:
  (i) the model quadratic $K (r^\ast)^2+2c\,r^\ast-1=0$, hence
  $(1-K (r^\ast)^2)/(2r^\ast)=c$;
  (ii) $0<r^\ast<1$;
  (iii) the \emph{bracket value}
  $c+K r^\ast=\sqrt{c^2+K}=:B>0$, the admissibility
  denominator evaluated at the level-$c$ model circle. For the sphere
  $B=\sqrt{1+c^2}\ge1$ is an absolute lower bound; for the hyperbolic plane
  $B=\sqrt{c^2-1}$ is positive but tends to $0$ as $c\to1^{+}$ (the
  escape-velocity boundary), so downstream margin constants scaled by $B$ are
  $c$-dependent, not absolute; for the flat plane $B=c$. Moreover the
  admissible windows themselves are captured by one identity
  (\texttt{centeredRadius\_lt\_one\_iff}): for any $c>0$,
  \[
    r^\ast(K,c)<1
    \iff
    c>\frac{1-K}{2},
  \]
  the uniform confinement threshold ($0$ at $S^2$, $1$ at $H^2$, $1/2$ flat)
  carried by \cref{def:spaceform_four_vertex}.
\end{lemma}

\begin{proof}
  \leanok
  All are elementary $\sqrt{\phantom{x}}$-algebra on the rationalized form:
  clear the denominator $\sqrt{c^2+K}+c>1$ (positive in all three
  admissible windows, \texttt{window\_pos}) and use
  $(\sqrt{c^2+K})^2=c^2+K$, then
  \texttt{linear\_combination}/\texttt{nlinarith}; the membership bounds
  compare $\sqrt{c^2+K}$ against $c$ and the threshold by strict
  monotonicity of the square root, and the threshold identity is a per-sign
  case split (at $K=0$, $r^\ast=1/(2c)<1\iff c>1/2$).
  \texttt{centeredRadius\_facts} packages (ii)--(iii) for the margins file.
\end{proof}

\begin{lemma}\leanok
[admissible uniform lower bound]
  \label{lem:spaceform_admissible_lower_bound}
  \lean{Gluck.SpaceForm.exists_admissible_lower_bound}
  \uses{def:spaceform_four_vertex, lem:sphere_curvature_lower_bound}
  From $\mathrm{SpaceFormFourVertex}\;K\;\kappa$ there is a
  confinement radius $R$ with $0<R<1$ and $R<\kappa(\theta)$ for all
  $\theta$. Together with \cref{lem:spaceform_denom_pos} this makes the
  admissibility denominator strictly positive on the slab
  $\{\|z\|\le R\}$ — the entry point of the whole flow layer.
\end{lemma}

\begin{proof}
  \leanok
  The minimum of the strictly positive continuous $\kappa$ over the compact
  period $[0,2\pi]$ is positive; take $R=\min(\min\kappa,1)/2$ and extend to
  all $\theta$ by $2\pi$-periodicity. ($K$-uniform: the hyperbolic
  bound $\kappa>1$ is not needed for \emph{this} lemma, only later for the
  margins scale via \cref{lem:centered_radius_spaceform_props}.)
\end{proof}

\subsection{Truncated flow}

\begin{definition}\leanok
[truncated gauge speed and reconstruction field]
  \label{def:spaceform_truncated_field}
  \lean{Gluck.SpaceForm.truncatedSpeed,
        Gluck.SpaceForm.truncatedField,
        Gluck.SpaceForm.truncatedSpeed_eq}
  \uses{def:spaceform_speed, def:truncated_speed, def:truncated_field}
  The \emph{truncated gauge speed}
  \[
    \hat q_{K,\kappa,R,\delta}(\theta,z)=
    \frac{1+K(\min(\|z\|,R))^2}
         {2\max\bigl(\kappa(\theta)-K\langle z,
           ie^{i\theta}\rangle,\ \delta\bigr)}
  \]
  and the \emph{truncated reconstruction field}
  $F_{K,\kappa,R,\delta}(\theta,z)=\hat q\cdot e^{i\theta}$, the
  right-hand side of the truncated reconstruction ODE $\gamma'=F(\theta,\gamma)$. On
  the admissible slab $\{\|z\|\le R,\ \delta\le\kappa-K\langle z,
  ie^{i\theta}\rangle\}$ both clamps are inactive and
  $\hat q=q_{K,\kappa}$ (\texttt{truncatedSpeed\_eq}). The clamped
  numerator stays $\ge1-R^2>0$ for $|K|\le1$ and $R<1$ (automatic
  in the disk), so the field is globally bounded by $(1+R^2)/(2\delta)$,
  Lipschitz in $z$ uniformly in $\theta$, and jointly continuous — in both
  models at once. $K$-generic transport of
  \cref{def:truncated_speed} and \cref{def:truncated_field}.
\end{definition}

\begin{definition}\leanok
[space-form flow and endpoint map]
  \label{def:spaceform_flow}
  \lean{Gluck.SpaceForm.spaceFormFlow, Gluck.SpaceForm.spaceFormEndpoint}
  \uses{def:spaceform_truncated_field, def:spherical_flow,
        def:spherical_endpoint}
  The \emph{Picard--Lindel\"of flow}
  $\Phi_{K,\kappa,R,\delta,r_0}:\mathbb C\times\mathbb R\to\mathbb
  C$ of the truncated field: a choice, made once per parameter tuple, of a
  map flowing every initial point of $\{\|z_0\|\le r_0\}$ along
  $F_{K,\kappa,R,\delta}$ on $[0,2\pi]$, jointly continuously
  (junk, $\mathrm{Prod.fst}$, when the side conditions $|K|\le1$,
  $\kappa$ continuous, $0\le R<1$, $0<\delta$ fail). The \emph{closing-error
  endpoint map} is $z_0\mapsto\Phi(z_0,2\pi)-z_0$; its interior zeros are the
  closed trajectories that the winding argument of
  \cref{chap:spaceform_converse} hunts. Transport of
  \cref{def:spherical_flow} and \cref{def:spherical_endpoint}.
\end{definition}

\begin{lemma}\leanok
[flow specification, uniqueness, and endpoint continuity]
  \label{lem:spaceform_flow_spec}
  \lean{Gluck.SpaceForm.spaceFormFlow_spec,
        Gluck.SpaceForm.spaceFormFlow_unique,
        Gluck.SpaceForm.spaceFormEndpoint_continuousOn,
        Gluck.SpaceForm.truncatedField_isPicardLindelof}
  \uses{def:spaceform_flow, def:spaceform_truncated_field,
        lem:spherical_flow_spec, lem:truncated_field_picard}
  For $\|z_0\|\le r_0$ the flow starts at $z_0$ and solves
  $\gamma'=F_{K,\kappa,R,\delta}(\theta,\gamma)$ on $[0,2\pi]$
  ($\mathrm{HasDerivWithinAt}$ form); any solution with the same initial
  value agrees with $\Phi(z_0,\cdot)$ on $[0,2\pi]$; and the endpoint map is
  continuous on $\{\|z_0\|\le r_0\}$.
\end{lemma}

\begin{proof}
  \leanok
  \uses{def:spaceform_truncated_field}
  One Picard--Lindel\"of package covers $[0,2\pi]$: the field is bounded by
  $B=(1+R^2)/(2\delta)$ and globally Lipschitz
  (\texttt{truncatedField\_isPicardLindelof}), so the budget condition is met
  by outer radius $r_0+2\pi B+1$, and Mathlib's
  \texttt{IsPicardLindelof} existence-with-continuous-dependence supplies the
  flow; the specification and continuity are
  \texttt{Classical.choose\_spec}. Uniqueness is
  \texttt{ODE\_solution\_unique\_of\_mem\_Icc\_right} against the uniform
  Lipschitz witness; endpoint continuity is the flow restricted to
  $\theta=2\pi$ minus the identity. Mirrors the landed proofs at
  \cref{lem:spherical_flow_spec} with $|K|\le1$ threaded through.
\end{proof}

\subsection{Invariant admissible domain}

\begin{lemma}\leanok
[invariant admissible domain]
  \label{lem:spaceform_invariant_admissible}
  \lean{Gluck.SpaceForm.invariant_admissible_domain,
        Gluck.SpaceForm.gronwall_L1_drive}
  \uses{def:spaceform_truncated_field, lem:spaceform_denom_pos,
        lem:invariant_admissible_domain, lem:gronwall_L1_drive}
  Gr\"onwall continuous dependence keeping a trajectory in the admissible
  slab: if a trajectory $\gamma$ of $F_{K,\kappa,R,\delta}$ starts close
  to a reference trajectory $\gamma_s$ (for curvature $\kappa'$) that stays
  $\mu$-interior to the slab ($\|\gamma_s\|\le R-\mu$ and
  $K\langle \gamma_s,ie^{i\theta}\rangle\le\kappa_0-\delta-\mu$ with
  $\kappa\ge\kappa_0$), and
  \[
    e^{2\pi L}\Bigl(\|\gamma(0)-\gamma_s(0)\|
      +\tfrac{1+R^2}{2\delta^2}\int_0^{2\pi}|\kappa-\kappa'|\Bigr)\le\mu,
  \]
  then $\gamma$ stays in the slab on $[0,2\pi]$: $\|\gamma(\theta)\|\le R$ and
  $\delta\le\kappa(\theta)-K\langle \gamma(\theta),
  ie^{i\theta}\rangle$. The drive is only small in $L^1$ (the Dahlberg
  regime), so the project-local Gr\"onwall lemma with $L^1$ drive is
  re-established here ($d(t)\le e^{LT}(d_0+\int_0^T g)$ from the integral
  inequality $d\le d_0+\int(Ld+g)$). Near-copy of
  \cref{lem:invariant_admissible_domain} with the $K$
  denominator; the hyperbolic numerator $1-\|z\|^2$ vanishes at the ideal
  boundary, so confinement is if anything easier than on the sphere.
\end{lemma}

\begin{proof}
  \leanok
  \uses{lem:spaceform_flow_spec}
  The curvature-sensitivity bound
  $|\hat q_{\kappa}-\hat q_{\kappa'}|\le M|\kappa-\kappa'|$ with
  $M=(1+R^2)/(2\delta^2)$ (\texttt{truncatedSpeed\_sub\_le}: shared clamped
  numerator in $[1-R^2,1+R^2]$, denominators $\ge2\delta$ differing by at
  most $2|\kappa-\kappa'|$ since $x\mapsto\max(x,\delta)$ is 1-Lipschitz)
  turns the trajectory difference into the integral inequality
  $\|\gamma-\gamma_s\|(t)\le\|\gamma(0)-\gamma_s(0)\|+\int_0^t(L\|\gamma-\gamma_s\|+M|\kappa-\kappa'|)$;
  Gr\"onwall with $L^1$ drive bounds $\|\gamma-\gamma_s\|\le\mu$ on $[0,2\pi]$, and the
  $\mu$-interior margins of $\gamma_s$ absorb the deviation
  (Cauchy--Schwarz for the inner-product margin, $|K|\le1$).
\end{proof}

\subsection{Arc algebra}

\begin{lemma}\leanok
[constant-curvature arcs are explicit circular arcs]
  \label{lem:spaceform_constant_curvature_arc}
  \lean{Gluck.SpaceForm.constant_curvature_arc,
        Gluck.SpaceForm.constant_arc_solves_truncated}
  \uses{def:spaceform_speed, def:spaceform_truncated_field,
        lem:constant_curvature_arc, lem:constant_arc_solves_truncated}
  Under the $K$-generic consistency identity
  $1+K\|w\|^2=2rk+K r^2$, at every angle $\theta$ where
  the bracket $k-K\langle \gamma(\theta), ie^{i\theta}\rangle$ stays
  positive, the circular arc $\gamma(\theta)=w-ire^{i\theta}$ has gauge speed
  exactly $r$ and solves the reconstruction ODE
  $\gamma'=q_k(\theta,\gamma)e^{i\theta}$; where moreover the clamps are inactive
  ($\|\gamma(\theta)\|\le R$, bracket $\ge\delta$) it solves the \emph{truncated}
  ODE for level $k$. Transport of \cref{lem:constant_curvature_arc} and
  \cref{lem:constant_arc_solves_truncated}: $K$ multiplies
  $\|w\|^2$ and $r^2$ in the consistency relation, not the $2rk$ term.
\end{lemma}

\begin{proof}
  \leanok
  The bracket identity $\langle \gamma(\theta), ie^{i\theta}\rangle=\langle w,
  ie^{i\theta}\rangle-r$ and the norm expansion
  $\|\gamma(\theta)\|^2=\|w\|^2-2r\langle w, ie^{i\theta}\rangle+r^2$ (both
  model-agnostic, copied verbatim from the Euclidean layer) reduce the speed
  claim to the consistency identity by field algebra; the derivative of
  $\theta\mapsto-ire^{i\theta}$ is $re^{i\theta}$, matching the field. On the
  admissible window \texttt{truncatedSpeed\_eq} removes the clamps.
\end{proof}

\begin{lemma}\leanok
[quadratic identity and sign-honest radius comparison]
  \label{lem:spaceform_speed_sub_radius}
  \lean{Gluck.SpaceForm.spaceFormSpeed_sub_radius,
        Gluck.SpaceForm.spaceFormSpeed_radius_le}
  \uses{def:spaceform_speed, def:centered_radius,
        lem:centered_radius_spaceform_props, lem:speed_quadratic_identity,
        lem:speed_radius_le}
  Exact second-order vanishing of the gauge speed at the centered circle:
  for the constant level $c$, $r^\ast=r^\ast(K,c)$, and any
  $(\theta,z)$ with $D=c-K\langle z, ie^{i\theta}\rangle\ne0$,
  \[
    q_c(\theta,z)-r^\ast
    \;=\;K\cdot
    \frac{\|z+r^\ast\, ie^{i\theta}\|^2}{2D}.
  \]
  Consequently $0\leK(q_c(\theta,z)-r^\ast)$ wherever $D>0$: for
  $K=+1$ model arcs stay \emph{outside} the centered circle
  ($r^\ast\le q_c$, as in \cref{lem:speed_radius_le}); for $K=-1$
  the inequality reverses ($q_c\le r^\ast$). The extra $K$ factor
  relative to the sphere identity \cref{lem:speed_quadratic_identity} is
  sign-critical: it is the ultimate source of the sign of the
  first-variation coefficient $\eta(K)$ in
  \cref{lem:spaceform_step_error_expansion}.
\end{lemma}

\begin{proof}
  \leanok
  Expand $\|z+r^\ast ie^{i\theta}\|^2=\|z\|^2+2r^\ast\langle z,
  ie^{i\theta}\rangle+(r^\ast)^2$ and eliminate $(r^\ast)^2$ with the model
  quadratic $K(r^\ast)^2+2cr^\ast-1=0$
  (\cref{lem:centered_radius_spaceform_props}); a
  \texttt{linear\_combination} closes the identity, and the comparison
  follows by multiplying through by $K$ ($K^2=1$).
\end{proof}

\begin{definition}\leanok
[arc-endpoint map]
  \label{def:spaceform_arc_map}
  \lean{Gluck.SpaceForm.spaceFormArcMap}
  \uses{def:spaceform_speed, def:spherical_arc_map}
  The endpoint of the constant-$k$ model arc of angular extent $\Delta$
  starting at $z$ with initial tangent angle $\theta_0$:
  \[
    A_{K,k,\theta_0,\Delta}(z)
    = z + i\,q_k(\theta_0,z)\,e^{i\theta_0}\,(1-e^{i\Delta}).
  \]
  Transport of \cref{def:spherical_arc_map}; it satisfies half-turn
  anti-equivariance
  $A_{K,k,\theta_0+\pi,\Delta}(-z)=-A_{K,k,\theta_0,\Delta}(z)$
  and concatenation
  $A_{\cdot,\Delta_2}\circ A_{\cdot,\Delta_1}=A_{\cdot,\Delta_1+\Delta_2}$
  along bracket-positive arcs
  (\texttt{spaceFormArcMap\_half\_turn}, \texttt{spaceFormArcMap\_concat}).
\end{definition}

\begin{definition}\leanok
[four-arc closing error map]
  \label{def:spaceform_step_error_map}
  \lean{Gluck.SpaceForm.stepErrorMap}
  \uses{def:spaceform_arc_map, def:step_error_map}
  The \emph{four-arc closing error map} $E^\ast_{K,a,b}$:
  $z+E^\ast_{K,a,b}(z)$ is the endpoint of the concatenated
  four-quarter-arc trajectory with levels $a,b,a,b$ at base angles
  $0,\pi/2,\pi,3\pi/2$ (the symmetric step model matching the canonical
  four-arc step curvature \cref{def:four_arc_step_function}). Transport of
  \cref{def:step_error_map}.
\end{definition}

\subsection{Margins and step-model transport}

\begin{definition}\leanok
[single-arc margin package]
  \label{def:spaceform_arc_margins}
  \lean{Gluck.SpaceForm.arcMargins}
  \uses{def:spaceform_speed, def:arc_margins}
  $\mathrm{arcMargins}\;K\;\kappa_0\;R\;\delta\;\mu\;K\;t_1\;t_2\;p$:
  along $[t_1,t_2]$ the constant-level-$k$ arc trajectory through $(t_1,p)$
  stays $\mu$-inside the norm clamp ($\|\cdot\|\le R-\mu$), $\mu$-inside the
  signed curvature margin
  ($K\langle\cdot, ie^{i\theta}\rangle\le\kappa_0-\delta-\mu$), and
  keeps the level-$k$ clamps inactive
  ($k-K\langle\cdot, ie^{i\theta}\rangle\ge\delta$). Transport of
  \cref{def:arc_margins}; exactly the hypothesis package consumed per
  quarter by \cref{lem:spaceform_step_model_transport}.
\end{definition}

\begin{lemma}\leanok
[uniform step-model margins near the centered circle]
  \label{lem:spaceform_step_model_margins}
  \lean{Gluck.SpaceForm.stepModel_margins}
  \uses{def:spaceform_arc_margins, def:spaceform_arc_map,
        def:centered_radius, lem:centered_radius_spaceform_props,
        lem:spaceform_speed_sub_radius, lem:step_model_margins}
  For $K\in\{+1,-1,0\}$, an admissible level $c$, and any
  $\kappa_0>-K r^\ast$ ($r^\ast=r^\ast(K,c)$), there are
  explicit $0<R<1$ and $\delta,\mu,\rho_0,h_0>0$ (functions of
  $c,\kappa_0,K$ only) such that for all levels $a,b$ within $h_0$
  of $c$ and every start $z_0$ within $\rho_0$ of the model-circle base point
  $z_0^\ast=-r^\ast i$, all four quarter-arc margin packages of the step
  model hold. Constants ledger, with $B=c+K
  r^\ast=\sqrt{c^2+K}$ and
  $m=\min(1-r^\ast,\ \kappa_0+K r^\ast,\ B)$: $R=(1+r^\ast)/2$,
  $\delta=\mu=m/8$, $\sigma=\min(r^\ast/2,m/32,B/4)$, $\rho_0=\sigma/32$,
  $h_0=\min(\sigma B^2/512,\,B/4)$. \textbf{$B$ — not $1$ — sets the
  scale}: this is the one place the sphere's absolute-constants shortcut
  (compare \cref{lem:step_model_margins}) fails for $K=-1$, where
  $B=\sqrt{c^2-1}\to0$ as $c\to1^{+}$.
\end{lemma}

\begin{proof}
  \leanok
  \uses{lem:spaceform_constant_curvature_arc}
  Quantitative perturbation off the centered circle: the gauge speed moves by
  at most $8h/B^2+d^2/B$ under a level shift $\le h$ and a start deviation
  $\le d$ (\texttt{spaceFormSpeed\_near\_circle}, driven by the exact level
  sensitivity and the quadratic identity of
  \cref{lem:spaceform_speed_sub_radius}); the whole arc then stays within
  $2d+16h/B^2$ of the centered circle (\texttt{arcDeviation\_bound}), and the
  deviation propagates across the four chained arcs with factor $2$ per arc,
  staying $\le\sigma$ throughout. A deviation $\le\sigma$ plus the three
  numeric slack inequalities yields each \texttt{arcMargins} package
  (\texttt{arcMargins\_of\_dev}): norm $\le r^\ast+\sigma\le R-\mu$, signed
  inner product $\le\sigma-K r^\ast\le\kappa_0-\delta-\mu$, level
  bracket $\ge(c-h_0)+K r^\ast-\sigma\ge\delta$.
\end{proof}

\begin{lemma}\leanok
[four-arc step-model transport]
  \label{lem:spaceform_step_model_transport}
  \lean{Gluck.SpaceForm.stepModel_transport,
        Gluck.SpaceForm.quarter_step_transport,
        Gluck.SpaceForm.invariant_admissible_arc}
  \uses{def:spaceform_arc_margins, def:spaceform_step_error_map,
        def:spaceform_truncated_field, def:four_arc_step_function,
        lem:spaceform_invariant_admissible,
        lem:spaceform_constant_curvature_arc,
        lem:step_model_transport, lem:quarter_step_transport,
        lem:invariant_admissible_arc}
  Compare a trajectory $\gamma$ of the $\kappa$-truncated flow on $[0,2\pi]$ with
  the concatenated four-quarter-arc step model from $\gamma_0$ (levels $a,b,a,b$,
  matching $\kappa^\ast=\mathtt{stepCurvature}\;b\;a\;0\;(\pi/2)\;\pi\;
  (3\pi/2)$, reused verbatim from \cref{def:four_arc_step_function}). If
  every quarter arc carries its \texttt{arcMargins} package and
  $e^{2\pi L}\,M\int_0^{2\pi}|\kappa-\kappa^\ast|\le\mu$
  ($M=(1+R^2)/(2\delta^2)$, $L$ a Lipschitz witness for the field), then $\gamma$
  is admissible on all of $[0,2\pi]$ and
  \[
    \bigl\|(\gamma(2\pi)-\gamma_0)-E^\ast_{K,a,b}(\gamma_0)\bigr\|
    \;\le\; e^{2\pi L}\,M\int_0^{2\pi}|\kappa-\kappa^\ast|.
  \]
  Transport of \cref{lem:step_model_transport} /
  \cref{lem:quarter_step_transport}; gains the hypotheses
  $|K|\le1$ and $R<1$ needed by the $K$-generic
  single-arc Gr\"onwall lemma.
\end{lemma}

\begin{proof}
  \leanok
  \uses{lem:spaceform_step_model_margins}
  Single quarter first (\texttt{quarter\_step\_transport}): the constant-$k$
  model arc solves the truncated ODE where its margins hold
  (\cref{lem:spaceform_constant_curvature_arc}), so the
  shifted-interval Gr\"onwall comparison
  (\texttt{invariant\_admissible\_arc}, the
  \cref{lem:spaceform_invariant_admissible} argument on $[t_1,t_2]$ against
  a constant-level reference, where the drive $M|\kappa-k|$ is continuous)
  gives admissibility on the quarter and an endpoint within the Gr\"onwall
  bound of the arc-map image $A_{K,k,t_1,t_2-t_1}(p)$. Then chain
  four quarters: the $L^1$ error splits over the quarters (the canonical
  step curvature is measurable and two-valued; the splitting lemmas are
  re-established locally since they are \texttt{private} in the base file),
  each quarter's start deviation is the previous quarter's endpoint bound,
  and the geometric accumulation
  $e^{2\pi L}(\dots)$ telescopes against the four-arc composite form of
  $E^\ast$ (\texttt{stepErrorMap\_four\_arc}).
\end{proof}

\subsection{First variation}

\begin{lemma}\leanok
[first-variation expansion of the step error map]
  \label{lem:spaceform_step_error_expansion}
  \lean{Gluck.SpaceForm.stepError_expansion}
  \uses{def:spaceform_step_error_map, def:centered_radius,
        lem:centered_radius_spaceform_props,
        lem:spaceform_speed_sub_radius, lem:step_error_expansion}
  The linchpin analytic estimate, transported from
  \cref{lem:step_error_expansion} (the highest-risk part of the layer: it is
  re-derived, not reused). For an admissible level $c$ there are
  $\rho_1,h_1,C>0$ such that for every step height $0<h\le h_1$ and every
  base point $z_0$ with $\|z_0+r^\ast i\|\le\rho_1$
  ($r^\ast=r^\ast(K,c)$),
  \[
    \Bigl\|E^\ast_{K,\,c-h/2,\,c+h/2}(z_0)
      +\eta(K)\,h\,\overline{(z_0+r^\ast i)}\Bigr\|
    \;\le\; C\,h\,\bigl(\|z_0+r^\ast i\|^2+h\bigr),
    \qquad
    \eta(K)=\frac{2\,K\, r^\ast}{c^2+K}.
  \]
  The linear-in-$h$ term is an anti-holomorphic conjugation with coefficient
  $\eta(K)$. The $K$ factor is mandatory — it descends
  from the $K$ on the right-hand side of
  \cref{lem:spaceform_speed_sub_radius} — so $\eta>0$ for the sphere,
  $\eta<0$ for the hyperbolic plane, and $\eta=0$ for the flat plane
  ($K=0$, where the four-arc error map is identically zero and the
  statement degenerates to the true bound
  $\|E^\ast\|\le Ch(\|\delta\|^2+h)$). The winding argument needs
  $\eta\ne0$ (\cref{lem:spaceform_step_error_coeff_ne_zero}) — which is
  where, and only where, the flat member leaves the endpoint-winding
  pipeline.
\end{lemma}

\begin{proof}
  \leanok
  Four bespoke per-arc estimates (\texttt{sf\_stepError\_arc1} through
  \texttt{arc3} plus the fourth inline): each quarter-arc speed is
  decomposed as
  $q=r^\ast+\eta$-linear term $+$ controlled remainder via the exact level
  sensitivity and the quadratic identity, with all error terms bounded by
  explicit fractions of $B=\sqrt{c^2+K}$ (radii and constants like
  $\min(1,c^2+K)/2^{26}$ and $C=1.6\times10^7(1/B+1/B^3)$ are
  fully explicit). The assembly identity collapses the four main terms to
  the single conjugation $-\eta h\overline{(z_0+r^\ast i)}$ and absorbs the
  residuals (an \texttt{nlinarith}-heavy final step under a raised
  heartbeat budget).
\end{proof}

\begin{lemma}\leanok
[nonvanishing first-variation coefficient]
  \label{lem:spaceform_step_error_coeff_ne_zero}
  \lean{Gluck.SpaceForm.stepError_coeff_ne_zero}
  \uses{def:centered_radius, lem:centered_radius_spaceform_props,
        lem:spaceform_step_error_expansion}
  For $K\in\{+1,-1\}$ (genuinely two-sign: the flat coefficient
  vanishes),
  $\eta(K)=2K r^\ast(K,c)/(c^2+K)
  \ne0$ for every admissible level $c$. Nonvanishing — \emph{not}
  positivity — is the property the winding argument of
  \cref{chap:spaceform_converse} consumes: the boundary loop of
  $E^\ast\approx-\eta h\,\overline{\delta}$ has winding number $-1$ for any
  real $\eta\ne0$, since scaling a loop by a nonzero real constant preserves
  its winding number. ($\eta>0$ at $K=+1$ but $\eta<0$ at
  $K=-1$.)
\end{lemma}

\begin{proof}
  \leanok
  $r^\ast>0$ by \cref{lem:centered_radius_spaceform_props}, $K\ne0$,
  and $c^2+K>0$ in both admissible branches; a quotient of
  nonzeros is nonzero.
\end{proof}
